在前面的讨论中，我们把气体看作大量自由运动的分子组成的系统，并用理想气体物态方程描述它们。
这一模型在气体密度较低、分子间距离较大时十分成功，但它有一个明显的局限：理想气体分子之间完全
没有相互作用，也没有自身大小，因此无论怎样降温、压缩，理想气体都不会自动变成液体或固体。真实物质
却可以在固、液、气三种物态之间转化。要理解这些相变，就必须把分子的有限大小和分子间作用力考虑进来。

从微观上看，物质的物态由两种趋势的竞争决定。热运动倾向于使分子分散、排列无序；分子间吸引作用
倾向于把分子聚集起来。当温度较低或密度较大时，分子间作用更重要，物质容易成为液体或固体；当温度
较高、密度较小时，热运动占优势，物质表现为气体。固体熔化成液体的过程叫熔化，液体凝固成固体的过程
叫凝固；液体变成气体叫汽化，气体变成液体叫液化；固体也可以不经过液态直接变成气态，这一过程叫升华。

\subsection*{相图}

物质在不同温度和压强下稳定存在的物态，可以用相图表示。图22.1是水的 $p$--$T$ 相图。图中三条曲线
分别是固液共存的熔化曲线、液气共存的汽化曲线和固气共存的升华曲线。曲线上每一点都表示两种物态能够
平衡共存；曲线把平面分为固态区、液态区和气态区。跨过一条相界线时，物质发生相变。

\phaseDiagramWater
\diagramcaption{图 22.1：水的 $p$--$T$ 相图及三种物态的微观示意。曲线、三相点和临界点的位置关系与水的相图特征一致。}

三条相界线交于点 $A$，称为三相点。在三相点，固态、液态和气态能够同时平衡共存。水的三相点温度为
$0.01\,{}^{\circ}\mathrm C$，压强为 $611.73\,\mathrm{Pa}$。液气共存曲线并不无限延伸，它终止于点 $C$，
称为临界点。水的临界温度约为 $374\,{}^{\circ}\mathrm C$，临界压强约为
$2.2\times10^7\,\mathrm{Pa}$。温度高于临界温度以后，不论怎样等温压缩，都不会出现清晰的气液分界面。

相变时系统需要吸收或放出热量。例如在一定压强下，固体达到熔点后继续吸热，温度可以保持不变而逐渐
熔化；液体达到沸点后继续吸热，温度也可以保持不变而逐渐汽化。单位质量物质完成某种相变所需的热量
称为相应的潜热。相变潜热说明：即使宏观温度不变，分子的空间结构和分子间势能仍会发生变化。

\subsubsection*{相变时为什么吸热而不升温}

对固定压强下的两相平衡，温度由共存条件确定。输入热量主要用于改变两相比例，
并不必然增大每个分子的平均平动能。设单位质量潜热为 $L$，转化质量为 $m$，则
$Q=mL$。但 $L$ 不是单位质量内能变化：若只有体积功，第一定律给出
\begin{equation}
 \Delta u=L-p(v_2-v_1),
\end{equation}
其中 $u$ 为比内能，$v$ 为比体积，单位分别为 $\mathrm{J/kg}$、$\mathrm{m^3/kg}$。
把 $h=u+pv$ 定义为比焓，在同一压强下就有 $L=h_2-h_1$。
气化时 $v_g-v_l>0$，吸热的一部分成为向外膨胀功；因此通常 $\Delta u<L$。
若过程包含显著电功、表面功或动能变化，就必须另列这些能量项，不能直接套用这一简式。

\subsubsection*{两相比例与杠杆关系}

在给定的气液共存温度下，令总质量为 $m$，其中气相质量为 $m_g$，气相质量分数
$x=m_g/m$。忽略界面体积，两相总体积满足
\[
 V=(m-m_g)v_l+m_gv_g,
 \qquad \bar v=\frac Vm=(1-x)v_l+xv_g.
\]
移项得
\begin{equation}
 x=\frac{\bar v-v_l}{v_g-v_l},\qquad 0\le x\le1.
\end{equation}
所以等温两相区的体积变化不是各相任意改变密度，而主要是气相、液相比例改变。
当 $\bar v=v_l$ 时全部为液体；当 $\bar v=v_g$ 时刚好全部成为蒸气。
这也解释了后面含水空气题中，液水尚存与液水恰好消失之后必须使用不同压强表达式。

\subsubsection*{共存曲线的斜率：克拉珀龙方程}

对固定组成、仅有体积功的纯物质，可逆基本关系为
$\dd U=T\dd S-p\dd V$。定义 $G=U+pV-TS$，逐项微分并约去对应项：
\[
 \dd G=(T\dd S-p\dd V)+p\dd V+V\dd p-T\dd S-S\dd T
 =V\dd p-S\dd T.
\]
对同一种物质的单位质量，$g=G/m$ 满足 $\dd g=v\dd p-s\dd T$。
两相平衡要求 $g_1=g_2$；沿共存曲线移动一小步后仍然相等，因此
\[
 v_1\dd p-s_1\dd T=v_2\dd p-s_2\dd T,
 \qquad (v_2-v_1)\dd p=(s_2-s_1)\dd T.
\]
相变可逆且等温时 $L=T(s_2-s_1)$，于是
\begin{equation}
 \frac{\dd p}{\dd T}=\frac{L}{T(v_2-v_1)}.
\end{equation}
两边的单位均为 $\mathrm{Pa/K}$。汽化时 $L>0,v_g>v_l$，共存曲线斜率为正；
普通条件下冰融化为水时体积反而减小，故水的固液共存线斜率为负。
这给出了相图中倾斜方向的物理依据，不应靠记忆图形判断。

在远离临界点的汽化过程中，若 $v_g\gg v_l$、蒸气近似理想气体，
$v_g=RT/(Mp)$，其中 $M$ 是摩尔质量，则
\[
 \frac{\dd\ln p_s}{\dd T}=\frac{ML}{RT^2}.
\]
再在不太宽的温区内把 $L$ 视为常数，从 $T_1$ 积分到 $T_2$ 得
\begin{equation}
 \ln\frac{p_s(T_2)}{p_s(T_1)}
 =\frac{ML}{R}\left(\frac1{T_1}-\frac1{T_2}\right).
\end{equation}
升温时右端为正，饱和蒸气压升高；临界点附近两相比体积接近，不能继续忽略 $v_l$。


\subsection*{Van der Waals 方程}

理想气体模型忽略了分子的大小和相互作用。Van der Waals 用两个修正把这两种效应纳入物态方程。第一，
分子本身占有体积，分子质心能够活动的有效体积小于容器的几何体积。若每个分子对应的排除体积参数为
$b$，有效体积写成 $V-Nb$。第二，容器内部的分子受到周围分子近似对称的吸引，而靠近器壁的分子受到
的吸引不对称，撞击器壁前会被向内拉回，因此实测压强小于没有吸引力时的压强。单位体积内分子数为
$N/V$，这一修正与分子对数成正比，量级为 $a(N/V)^2$。于是得到
\begin{equation}
  \boxed{\left[p+a\left(\frac{N}{V}\right)^2\right](V-Nb)=Nk_B T}.
  \tag{22.1}
\end{equation}
这里 $a$ 衡量分子间吸引作用，$b$ 衡量分子的有限大小。若以物质的量 $n$ 表示，同一个方程也常写为
\[
  \left(p+a_m\frac{n^2}{V^2}\right)(V-nb_m)=nRT,
\]
此时 $a_m=N_A^2a$、$b_m=N_Ab$ 是摩尔参数，单位分别为
$\mathrm{Pa\,m^6\,mol^{-2}}$ 和 $\mathrm{m^3\,mol^{-1}}$，$N_A$ 为阿伏伽德罗常量。
表22.1中的 $a,b$ 按传统表头表示这两个摩尔参数，不能与分子形式的数值混用。
排除体积参数也不是分子自身几何体积：硬球模型中单粒子参数是球自身体积的4倍，来自两球球心的不可重叠区域及成对计数。
表22.1列出若干物质的经验参数；不同物质的参数不同，反映了分子大小和吸引作用
强弱的差异。

\begin{table}[htbp]
  \centering
  \caption{表 22.1：若干物质的 Van der Waals 参数}
  \renewcommand{\arraystretch}{1.12}
  \begin{tabular}{ccc}
    \toprule
    物质 & $a/(\mathrm{Pa\,m^6\,mol^{-2}})$ & $b/(10^{-3}\,\mathrm{m^3\,mol^{-1}})$ \\
    \midrule
    $\mathrm H_2$   & 0.01945 & 0.022 \\
    $\mathrm H_2O$ & 0.56539 & 0.031 \\
    $\mathrm N_2$  & 0.13882 & 0.039 \\
    $\mathrm O_2$  & 0.13983 & 0.032 \\
    $\mathrm{CO}_2$& 0.37186 & 0.043 \\
    \bottomrule
  \end{tabular}
\end{table}

Van der Waals 方程仍然是近似模型，但它已经能同时描述稀薄气体、液体的有限体积以及气液相变的主要
特征。它的重要性不只在于提高数值精度，更在于用一个简单方程把气态和液态放进同一幅物理图景中。

\subsection*{气液相变}

把式(22.1)改写为
\[
  p=\frac{Nk_BT}{V-Nb}-a\left(\frac NV\right)^2.
\]
在密度很低、$Nb/V\ll1$ 时，可以把第一项按 $Nb/V$ 展开：
\begin{equation}
  \boxed{p=\frac{Nk_BT}{V}
  +\left(\frac NV\right)^2(k_BT b-a)+\cdots }.
  \tag{22.2}
\end{equation}
第一项就是理想气体压强；第二项同时包含有限分子体积造成的正修正和分子吸引造成的负修正。当温度较高
时，$k_BT b$ 项更重要，有限体积使压强偏大；温度较低时，$a$ 项可能占优势，吸引作用使压强偏小。

图22.2给出了 Van der Waals 方程在不同温度下的等温线。温度高于某一特征温度时，等温线随体积增加
单调下降，与普通气体相似；温度降低后，等温线上出现一段波状回环。临界温度 $T_c$ 恰好把这两种情形
分开，临界等温线在点 $C$ 处出现水平拐点。图中从下到上的温度依次为
$T_1=0.85T_c,T_2=0.90T_c,T_c,T_3=1.08T_c,T_4=1.16T_c,T_5=1.24T_c$。
图中先用无量纲比值 $t'=T/T_c$ 标出这六个温度；坐标 $v'=V/V_c$、$p'=p/p_c$
也分别用临界值作单位，临界参数的求法将在下一节推导。

\vdWIsotherms
\diagramcaption{图 22.2：Van der Waals 方程的等温线。低于临界温度时出现回环，$T_c$ 等温线在临界点出现水平拐点。}

方程给出的整段回环并不全是均匀物质可以稳定经历的状态。真实物质在低于临界温度时等温压缩，达到
饱和蒸气状态后会开始液化；在一段有限体积范围内，液体和蒸气共同存在，压强保持在该温度的饱和蒸气压。
随着体积继续减小，蒸气所占比例下降、液体所占比例上升，直到物质全部变为液体。如何用理论等温线确定
这段恒压相变线，将在后面的 Maxwell 等面积法中讨论。

\subsection*{临界点}

临界点是气液共存曲线的终点。图22.3(a)中，临界等温线在 $C$ 点具有水平拐点，所以
\[
  \left(\frac{\partial p}{\partial V}\right)_{T_c}=0,
  \qquad
  \left(\frac{\partial^2 p}{\partial V^2}\right)_{T_c}=0.
\]
把 Van der Waals 方程逐次求导，两个条件为
\[ -\frac{Nk_BT_c}{(V_c-Nb)^2}+\frac{2aN^2}{V_c^3}=0,
\qquad \frac{2Nk_BT_c}{(V_c-Nb)^3}-\frac{6aN^2}{V_c^4}=0.\]
分别移项后相除，得 $2/(V_c-Nb)=3/V_c$，因此 $V_c=3Nb$。
代回第一式，$Nk_BT_c=2aN^2(2Nb)^2/(3Nb)^3=8aN/(27b)$；
最后代回状态方程即可求出临界压强。于是
\begin{equation}
  \boxed{p_c=\frac{a}{27b^2},\qquad
  T_c=\frac{8a}{27k_Bb},\qquad V_c=3Nb}.
  \tag{22.3}
\end{equation}

\Needspace{18\baselineskip}
\criticalPointFigure
\diagramcaption{图 22.3：临界点的两种表示。(a) 临界等温线的水平拐点；(b) $p$--$T$ 相图中气液共存线止于临界点。}

图22.3(b)区分了两条路径：水平的 $P\to M\to Q$ 在 $M$ 穿过气液共存线，会发生相变；
另一条 $P\to A\to B\to Q$ 先增压，再进入高于临界温度的区域，绕过临界点而不穿过共存线。
沿后一条路径，物质的密度可以连续变化，途中不会出现突变
的气液界面。由此可见，气体和液体并不是两种在任何条件下都截然分开的物态；在临界点以上，二者可以
连续过渡。

在透明高压容器中观察临界过程，会看到图22.4所示的现象。低于临界温度时，容器内有清晰的液面；接近
临界温度时，液体与蒸气的密度趋于相同，密度涨落对光产生强烈散射，系统呈现乳白色浑浊，称为临界乳光；
继续升温，使温度和压强都超过临界值后，界面消失，容器中成为均匀的超临界流体。
仅有 $T>T_c$ 并不足以称为超临界流体，还要有 $p>p_c$；若温度虽高而压强很低，仍是低密度气态。
下图表示选定充装量下经过临界区域的三个阶段，
使界面、密度涨落和均匀状态可以直接比较。

\criticalApparatusFigure
\diagramcaption{图 22.4：临界点附近的观察情景与状态示意。左：清晰气液界面；中：临界乳光；右：均匀超临界流体。}

临界点附近还有一个重要特征：等温压缩系数
\[
  \kappa_T=-\frac1V\left(\frac{\partial V}{\partial p}\right)_T
\]
变得很大，并在理想临界模型中发散。这意味着很小的压强变化就能引起很大的密度变化，也解释了临界点
附近强烈的密度涨落。

把压强、体积和温度分别用临界值作单位，定义约化变量
\begin{equation}
  \boxed{p'=\frac{p}{p_c},\qquad v'=\frac{V}{V_c},\qquad t'=\frac{T}{T_c}}.
  \tag{22.4}
\end{equation}
代入时，$p+aN^2/V^2=[a/(27b^2)](p'+3/v'^2)$，
$V-Nb=Nb(3v'-1)$，$Nk_BT=8aNt'/(27b)$。
两边约去共同因子 $aN/(27b)$，参数 $a,b$ 完全消失，得到
\begin{equation}
  \boxed{\left(p'+\frac{3}{v'^2}\right)(3v'-1)=8t'}.
  \tag{22.5}
\end{equation}
这说明在 Van der Waals 模型中，不同物质只要用各自临界参数作单位，就具有同一形式的物态方程和相似
的相变行为，这一结论称为对应态定律。真实物质不完全服从简单的 Van der Waals 方程，但约化变量和对应
态思想仍然非常有用。

\subsection*{低于临界温度的等温线}

在图22.5中，蓝色曲线是某一 $T<T_c$ 时 Van der Waals 方程给出的等温线。真实物质从体积较小的液态
$M$ 出发等温膨胀时，不会沿曲线走过 $M\to A\to P\to B\to N$ 的全部波状部分，而是沿水平直线
$M\to P\to N$ 发生气液共存相变。$M$ 是饱和液体，$N$ 是饱和蒸气，水平线压强 $p_T$ 就是温度 $T$
下的饱和蒸气压。

\maxwellAreaFigure
\diagramcaption{图 22.5：Maxwell 等面积作图。水平相变线的位置由上下两块面积相等确定。}

水平线的位置不能任意选择。设饱和液体和饱和蒸气的体积分别为 $V_l,V_g$。
令亥姆霍兹自由能 $F=U-TS$，在固定温度和粒子数下 $\dd F=-p\dd V$。
对模型的均匀态自由能作解析延拓，得到
\[F(V_g)-F(V_l)=-\int_{V_l}^{V_g}p(V)\dd V.\]
延拓曲线中的不稳定段不是实际可走过的稳定准静态过程。
平衡共存要求两相的吉布斯自由能 $G=F+p_TV$ 相等（比较同样粒子数），
即 $F(V_g)-F(V_l)=-p_T(V_g-V_l)$。二式比较就得到等面积关系。
也可从内能是状态函数理解这个形式积分：模型延拓给出
\[
  \Delta_1U=U(N)-U(M)
  =T\,[S(N)-S(M)]-\int_{V_l}^{V_g}p(V)\,\dd V.
\]
若沿真实的恒压相变线从 $M$ 到 $N$，同一对初、末态的内能改变量为
\[
  \Delta_2U=T\,[S(N)-S(M)]-p_T(V_g-V_l).
\]
内能是状态函数，$\Delta_1U=\Delta_2U$，所以
\begin{equation}
  \boxed{\int_{V_l}^{V_g}\bigl[p(V)-p_T\bigr]\,\dd V=0}.
  \tag{22.6}
\end{equation}
也就是说，水平线上方的面积 $S_2$ 必须等于水平线下方的面积 $S_1$。这就是 Maxwell 等面积法。它把
理论方程中不稳定的回环部分替换为真实的恒压气液共存线，同时保证相变两端具有相同的热力学势。

\subsection*{过冷气体与过热液体}

虽然平衡态相变线由 Maxwell 等面积法确定，但实际系统不一定恰好在到达饱和状态时立即相变。液滴的
形成需要一个凝结核，气泡的形成也需要一个成核中心；如果容器非常洁净，系统中缺少尘埃、划痕、自由
表面等成核条件，蒸气可以越过饱和点而暂时保持气态，液体也可以越过沸点而暂时保持液态。这些状态分别
称为过冷蒸气（或过饱和蒸气）和过热液体，它们是亚稳态。
在图22.5中，$M\to A$ 是过热液体分支，$N\to B$ 是过饱和蒸气分支。
它们尚满足 $(\partial p/\partial V)_T<0$，即等温压缩系数为正；但相分离后的自由能更低。
极值点 $A,B$ 是失稳界限（spinodal）；其间 $A\to P\to B$ 有
$(\partial p/\partial V)_T>0$，压缩系数为负，连局部力学稳定性也不满足。

亚稳态并非真正稳定。一旦出现凝结核或气泡核，相变会迅速发生。过冷蒸气遇到尘埃、离子或飞机尾气中
的微粒，就可能凝结成小液滴或冰晶。飞机在高空留下的白色尾迹，常常就是尾气中的水蒸气在低温空气中
凝结、冻结形成的。云的形成也与空气上升冷却并达到水蒸气饱和有关。

图22.6中，尾迹从发动机后方开始，随混合扩散逐渐展宽。它并非水蒸气本身的颜色，而是凝结、冻结后的微小颗粒散射光所形成的可见轨迹。
\contrailFigure
\diagramcaption{图 22.6：飞机尾迹的形成情景。可见尾迹来自水蒸气在微粒上凝结或凝华。}

亚稳态对粒子探测也十分重要。云室利用过饱和蒸气：带电粒子通过时沿途产生离子，水滴在离子上凝结，
于是粒子轨迹变成可见雾迹。气泡室则利用过热液体：带电粒子经过时触发微小气泡，形成可见径迹。云室
和气泡室曾经是发现基本粒子的重要工具，与这些技术有关的工作获得过诺贝尔物理学奖。它们共同利用了
同一原理：把系统准备在接近失稳的亚稳态，让微小扰动被迅速放大成可见的宏观结构。

\subsection*{表面张力、附加压强与润湿}

\subsubsection*{先数界面，再列受力}

把液体表面面积记为 $A$。在温度和组成给定时，缓慢增加面积所需的可逆功为
$\delta W_{\rm on}=\sigma\dd A$，$\sigma$ 的单位既可写成 $\mathrm{J/m^2}$，
也可写成 $\mathrm{N/m}$。面积功与表面内能并不总是同一件事：若 $\sigma$ 随温度改变，
拉伸时还可能吸、放热，相关熵项已在第二定律的表面工质题中讨论。

考虑半径 $R$ 的球形液滴，或忽略膜厚的球形肥皂泡。沿赤道切开，取上半部作为受力对象。
表面张力在赤道边缘沿球面切线作用，其竖直分量向下。图\ref{s31:fig:laplace-cut}
用竖直剖面标明切口；三个向上的箭头表示压差在水平投影面上的合力方向，
并不表示压强只作用在三个点上。

\begin{center}
\begin{tikzpicture}[x=1cm,y=1cm,>=Latex,font=\small]
 \draw[Ink,thick] (-1.5,0) arc[start angle=180,end angle=0,radius=1.5];
 \draw[TextGray,dashed] (-1.5,0)--(1.5,0);
 \fill[Ink] (0,0) circle (1pt);
 \draw[<->,TextGray] (0,-.25)--(1.5,-.25);
 \node[below] at (.75,-.25) {$R$};
 \foreach \x in {-.75,0,.75}{\draw[->,DeepBlue,thick] (\x,.10)--(\x,.82);}
 \draw[->,DeepBlue,thick] (-1.5,0)--(-1.5,-.72);
 \draw[->,DeepBlue,thick] (1.5,0)--(1.5,-.72);
 \node[anchor=east] at (-1.72,-.35) {边缘张力};
 \node[anchor=west] at (1.75,.75) {$\Delta p\,\pi R^2$ 向上};
 \node[anchor=west] at (1.75,-.35) {$k\sigma\,2\pi R$ 向下};
 \node at (0,1.85) {上半球受力的竖直剖面};
\end{tikzpicture}
\diagramcaption{球面沿赤道截开。$k$ 是界面数；每个界面在赤道各处都贡献向下的切向拉力。}\label{s31:fig:laplace-cut}
\end{center}

压差的合力等于 $\Delta p\,\pi R^2$，不是压差乘以半球曲面面积。
若有 $k$ 个具有相同张力系数、半径近似相同的界面，边缘张力合力为 $k\sigma\,2\pi R$。
平衡条件给出
\begin{equation}
 \Delta p\,\pi R^2=k\sigma\,2\pi R
 \quad\Longrightarrow\quad
 \Delta p=\frac{2k\sigma}{R}.
\end{equation}
液滴只有一个液气界面，$k=1$；薄肥皂泡有内外两个液气界面，$k=2$，故分别得到
$2\sigma/R$ 和 $4\sigma/R$。液体内部单个气泡通常也只有一个气液界面，不能因名称中有“泡”
便取 $k=2$。这些公式忽略重力造成的明显形变；$R$ 增大时附加压强减小，平面极限趋于零。

\subsubsection*{接触角如何决定液面升降}

接触角 $\theta$ 在液体一侧量取，是固体壁面与液气界面切平面之间的夹角。
圆管内半径为 $r$，液体密度为 $\rho$，两侧上方气压相同。管壁一周的表面张力合力
竖直分量为 $2\pi r\sigma\cos\theta$。当管径足够小、液面弯曲部分的体积可相对
$\pi r^2h$ 忽略时，与高出外液面的液柱重力平衡：
\[
 2\pi r\sigma\cos\theta=\rho g\pi r^2h,
 \qquad h=\frac{2\sigma\cos\theta}{\rho gr}.
\]
$h$ 是有符号的液面高度差。润湿时 $\theta<90^\circ$，$h>0$，液面升高；
不润湿时 $\theta>90^\circ$，$h<0$，液面降低。若先假定表面张力必向上，再把负号丢掉，
就会把凸液面画成错误的上升液柱。接触角、凹凸形状与液面升降必须同时核对。

